% concatenated from alpha-convex.tex
\documentclass[reqno]{amsart}
\usepackage[utf8]{inputenc}
\usepackage[a4paper]{geometry}

\usepackage[english]{babel}
\usepackage[document]{ragged2e}
\usepackage{amsfonts}
\usepackage{amsthm}
\usepackage{amssymb, amsmath}
\usepackage{hyperref}

\usepackage{xcolor}
\usepackage{cleveref}
\usepackage{enumitem}
\usepackage{graphicx}

\theoremstyle{plain}
\newtheorem{theo}{Theorem}[section]
\newtheorem{prop}[theo]{Proposition}
\newtheorem{lemma}[theo]{Lemma}
\newtheorem{coro}[theo]{Corollary}
\newtheorem{defi}{Definition}
\newtheorem{question}{Question}
\theoremstyle{definition}
\newtheorem{rem}[theo]{Remark}
\newtheorem{example}[theo]{Example}

\newcommand{\C}{\mathbb{C}}
\newcommand{\D}{\mathbb{D}}
\newcommand{\R}{\mathbb{R}}
\newcommand{\Z}{\mathbb{Z}}

\title[$\alpha$-Convex Functions]{A Characterization of $\alpha$-Convex Functions with Sharp Coefficient and Schwarzian Estimates}

\date{\today}

\author[V. Bravo]{V\'ictor Bravo}
\address{V\'ictor Bravo, Facultad de Ingenier\'ia y Ciencias, Universidad Adolfo Ib\'a\~nez, Av. Padre Hurtado 750, Vi\~na del Mar, Chile}
\email{victor.bravo.g@uai.cl}

\author[P. Carrasco]{Pablo Carrasco}
\address{Pablo Carrasco, Facultad de Ingenier\'ia y Ciencias, Universidad Adolfo Ib\'a\~nez, Av. Padre Hurtado 750, Vi\~na del Mar, Chile}
\email{pablo.carrasco.r@uai.cl}

\author[R. Hern\'andez]{Rodrigo Hern\'andez}
\address{Rodrigo Hern\'andez, Facultad de Ingenier\'ia y Ciencias, Universidad Adolfo Ib\'a\~nez, Av. Padre Hurtado 750, Vi\~na del Mar, Chile}
\email{rodrigo.hernandez@uai.cl}

\author[O. Venegas]{Osvaldo Venegas}
\address{Osvaldo Venegas, Departamento de Ciencias Matem\'aticas y F\'isicas, Facultad de Ingenier\'ia, Universidad Cat\'olica de Temuco, Av. Rudecindo Ortega 02950, Temuco, Chile}
\email{ovenegas@uct.cl}

\begin{document}

\maketitle

\justifying

\begin{abstract}
The class $M_\alpha$ of $\alpha$-convex functions, introduced by Mocanu in 1969, interpolates between starlike and convex functions. We prove a characterization of $M_\alpha$ that extends a theorem of Chuaqui, Duren, and Osgood from the convex case to the full class, and determine sharp values of $\beta$ for which $M_\alpha \subset C_\beta$ and $C_\beta \subset M_\alpha$. We also obtain a sharp Fekete--Szeg\H{o} inequality, bounds for the order and the Schwarzian norm, and an explicit formula for the Schwarzian norm of the $\alpha$-Koebe function for $\alpha = 1/n$, $n \in \mathbb{N}$, which we verify for $n \leq 9$ and conjecture to hold in general.
\end{abstract}
\vspace{0.3cm}
\textbf{Key words:} $\alpha$-convex functions, univalent functions, Schwarzian derivative. \\

\textbf{Mathematics subject classification:} 30C45, 30C55, 30C80. \\

\section{Introduction}\label{sec1}

The class $M_\alpha$ of \textit{$\alpha$-convex functions} was introduced by Mocanu \cite{Moc1969} in 1969 as a generalization of convex functions within the theory of univalent functions that interpolate continuously between starlike and convex mappings, controlled by the real parameter $\alpha$. A function $f(z) = z + a_2z^2 + \cdots$, analytic in the unit disk $\mathbb{D} = \{z : |z| < 1\}$, belongs to $M_\alpha$ if
\begin{equation}\label{eq:a-convex}
\operatorname{Re}\{J_\alpha[f](z)\} = \operatorname{Re}\left\{(1-\alpha)\frac{zf'(z)}{f(z)} + \alpha\left(1 + \frac{zf''(z)}{f'(z)}\right)\right\} > 0, \quad z \in \mathbb{D},
\end{equation}
where $\alpha \in \mathbb{R}$. In particular, $\alpha = 0$ recovers starlikeness and $\alpha = 1$ convexity. Mocanu proved that $M_\alpha \subseteq C$ for $\alpha \geq 1$, and Miller, Mocanu, and Reade \cite{MMR1978} proved that $M_\alpha \subset S^*$ for every $\alpha \in \mathbb{R}$, where $C$ and $S^*$ denote the classical classes of
convex and starlike functions. Moreover, for $\alpha \geq 1$ the same authors showed that $M_\alpha \subseteq S^*_\beta$, where
\[\beta = \frac{\Gamma\!\left(\frac{1}{2}+\frac{1}{\alpha}\right)}{\sqrt{\pi}\cdot
\Gamma\!\left(1+\frac{1}{\alpha}\right)},\]
and $S^*_\beta$ denotes the class of starlike functions of order $\beta$,
defined by
\[\operatorname{Re}\left\{J_0[f](z)\right\} = \operatorname{Re}\left\{\frac{zf'(z)}{f(z)}\right\} > \beta, \quad z \in \mathbb{D}.\]
The bound is sharp, with equality for the $\alpha$-Koebe function
\begin{equation}\label{koebeconvex}
k_\alpha(z) = \left[\frac{1}{\alpha}\int_0^z \xi^{1/\alpha-1}(1-\xi)^{-2/\alpha}\,
d\xi\right]^{\alpha}.
\end{equation}
In fact, the authors in \cite{MMR1974} show that $f\in M_\alpha$ ($\alpha>0$) if and only if there exists $g\in S^*$ such that
\begin{equation}\label{rep_int}
    f(z)=\left[\dfrac{1}{\alpha}\int_0^z\dfrac{(g(\xi))^{1/\alpha}}{\xi}\,d\xi\right]^\alpha.
\end{equation}
The previous example holds when $g(z)$ is the Koebe functions $k(z)=z/(1-z)^2$.

It is therefore natural to ask whether a more precise statement can be made on the convex side: for $\alpha \geq 1$, does every $f \in M_\alpha$ belong to $C_\beta$ for some $\beta > 0$? We shall prove that the answer is yes, with $\beta = \beta(\alpha) \in [0,1)$ explicit and sharp, where $C_\beta$ denotes the class of convex functions of order $\beta$, characterized
by
\[\operatorname{Re}\left\{J_1[f](z)\right\} = \operatorname{Re}\left\{1 +
\frac{zf''(z)}{f'(z)}\right\} > \beta, \quad z \in \mathbb{D}.\]
Propositions~\ref{prop_beta} and~\ref{prop_alpha} together make the reciprocal connection between $M_\alpha$ and $C_\beta$ precise.

The main result of this paper is the following characterization of $M_\alpha$, which
generalizes a theorem of Chuaqui, Duren, and Osgood \cite{CDO2011} to this setting.

\begin{theo}\label{thm:main}
Let $\alpha > 0$. Then $f \in M_\alpha$ if and only if
\begin{equation}\label{main}
\operatorname{Re}\{J_\alpha[f](z)\} \geq \frac{1}{4}(1 - |z|^2) \left|
(1-\alpha)\frac{zf'(z) - f(z)}{zf(z)} + \alpha\frac{f''(z)}{f'(z)}
\right|^2, \quad z \in \mathbb{D}.
\end{equation}
\end{theo}

The class $M_\alpha$ is not linearly invariant, which makes the usual invariant quantities, such as the order or the norm of the Schwarzian derivative, difficult to compute or estimate. In this paper, we give a complete determination of the order of $M_\alpha$ in terms of $\alpha$, and a partial answer for the Schwarzian norm. Concretely, we establish a sharp coefficient estimate for $f(z) = z + a_2 z^2 +
a_3 z^3 + \cdots$ in $M_\alpha$:
\[|a_3 - a_2^2| \leq \frac{1}{1 + 2\alpha}\left[1 + \left(\frac{|1-\alpha^2| - |1+\alpha|^2}{4}\right)|a_2|^2\right],\]
which refines existing bounds and highlights the influence of $\alpha$ on the coefficients of functions in $M_\alpha$. When $\alpha = 1$, the inequality reduces to the classical bound for $C$. In the final subsection, we compute the Schwarzian norm of the $\alpha$-Koebe function given by \eqref{koebeconvex} for $\alpha = 1/n$, $n \in \mathbb{N}$, and formulate a conjecture for the full class $M_\alpha$.

\section{Results}\label{sec2}

Let $f$ be holomorphic in $\mathbb{D}$, normalized by $f(0) = 0$ and $f'(0) = 1$, with $f(z)f'(z)/z \neq 0$ for all $z \in \mathbb{D}$. For $\alpha \in \mathbb{R}$, we say that $f$ is $\alpha$-convex if it satisfies \eqref{eq:a-convex}, and we write $f \in M_\alpha$. Note that $M_0 = S^*$ and $M_1 = C$. By \cite{MMR1973}, every
function in $M_\alpha$ is univalent and starlike, and $M_\alpha \subseteq C$ whenever $\alpha \geq 1$.

The classes $S^*_\beta$ and $C_\beta$ of starlike and convex functions of order $\beta$ ($0 \leq \beta < 1$) are classical extensions of $S^*$ and $C$. They satisfy $C_\beta \subset S^*_\gamma$, where
\begin{equation}\label{eq:gamma}
\gamma = \gamma(\beta) = \begin{cases}
\dfrac{1-2\beta}{2(2^{1-2\beta}-1)} & \text{if } \beta \neq \dfrac{1}{2}, \\[6pt]
\dfrac{1}{2\ln 2} & \text{if } \beta = \dfrac{1}{2},
\end{cases}
\end{equation}
see \cite{MM2000,GK2003}. When $\beta = 0$ this reduces to the classical inclusion $C \subset S^*_{1/2}$. Notice that from \eqref{eq:a-convex} and the fact that
\( M_\alpha\subseteq S^*_\delta \) (see \cite{MMR1978}) with
\[\delta=\dfrac{\Gamma\left(\frac{1}{2}+\frac{1}{\alpha}\right)}{\sqrt{\pi}\cdot\Gamma\left(1+\frac{1}{\alpha}\right)},\]
we have that
\begin{equation*}
    \text{Re}\left\{1+z\dfrac{f''}{f'}(z)\right\}>\left(\dfrac{\alpha-1}{\alpha}\right)\text{Re}\left\{z\dfrac{f'}{f}(z)\right\}>\left(\dfrac{\alpha-1}{\alpha}\right)\cdot \dfrac{\Gamma\left(\frac{1}{2}+\frac{1}{\alpha}\right)}{\sqrt{\pi}\cdot\Gamma\left(1+\frac{1}{\alpha}\right)}.
\end{equation*}
We have proved the following proposition.
\begin{prop}\label{prop_beta}
    Let $f\in M_\alpha$ with $\alpha\geq 1$, then $f\in C_{\beta}$, where
    \begin{equation*}
     \beta=\left(\dfrac{\alpha-1}{\alpha}\right)\cdot \dfrac{\Gamma\left(\frac{1}{2}+\frac{1}{\alpha}\right)}{\sqrt{\pi}\cdot\Gamma\left(1+\frac{1}{\alpha}\right)}.
    \end{equation*}
\end{prop}
The fact that \( M_\alpha \subseteq C \) for all \( \alpha\geq 1\) and that \( C_\beta\subset C=M_1 \) for all \( 0\leq \beta<1 \) allows us to formulate the next question: given \( f\in C_\beta \), does exist \( \alpha=\alpha(\beta) \) such that \( f\in M_\alpha\)? The next proposition gives a positive answer as a direct consequence of \eqref{eq:gamma}.

\begin{prop}\label{prop_alpha}
Let \( f\in C_\beta \) with \(0 \leq \beta <1 \), and  \( \gamma \) is given in \eqref{eq:gamma}. Then \( f\in M_\alpha \) with
\[\alpha=1+\dfrac{\beta}{\gamma-\beta}.\]
\end{prop}
\begin{proof}
    Let \(0\leq \beta <1 \). We know that if \( f\in C_\beta\) then
    \[\text{Re}\left\{1+z\dfrac{f''}{f'}(z)\right\}>\beta.\]
    By \eqref{eq:gamma} we conclude that \( f\in S^*_\gamma \) so that
    \[\text{Re}\left\{z\dfrac{f'}{f}(z)\right\}>\gamma. \]
    If \( \alpha=1+\beta/(\gamma-\beta) \) then
    \begin{align*}
      \text{Re}\left\{(1-\alpha)z\dfrac{f'}{f}(z)+\alpha\left(1+z\dfrac{f''}{f'}(z)\right)\right\}& > (1-\alpha)\gamma+\alpha\beta\\
      & =\left(1-\left(1+\dfrac{\beta}{\gamma-\beta}\right)\right)\gamma+\left(1+\dfrac{\beta}{\gamma-\beta}\right)\beta\\
      & = \dfrac{-\beta\gamma}{\gamma-\beta}+\dfrac{\beta\gamma}{\gamma-\beta}=0.
    \end{align*}
Then the proof is finished.
\end{proof}

Now we give the proof of Theorem~\ref{thm:main}.
\begin{proof}[\textbf{Proof of Theorem~\ref{thm:main}}]
If $f$ satisfies \eqref{main}, then in particular $\operatorname{Re}\{J_\alpha[f](z)\} \geq 0$, so $f \in M_\alpha$. Conversely, suppose $f \in M_\alpha$. Then $J_\alpha[f](z)$ belongs to the Carath\'eodory class $P$ of holomorphic functions with positive real part normalized by $J_\alpha[f](0) = 1$. Every function in $P$
is subordinate to $(1+z)/(1-z)$, so there exists a holomorphic
$\varepsilon \colon \mathbb{D} \to \mathbb{D}$ with $\varepsilon(0) = 0$ such that
\begin{equation}\label{eq:J_subord}
J_\alpha[f]=(1-\alpha)z\dfrac{f'}{f}(z)+\alpha\left(1+z\dfrac{f''}{f'}(z)\right)=\dfrac{1+\varepsilon(z)}{1-\varepsilon(z)},\ z\in \D.
\end{equation}
Subtracting $1$ from both sides of equation \eqref{eq:J_subord} and considering that, by Schwarz lemma, $\varepsilon(z)=z\varepsilon_0(z)$, for some $\varepsilon_0$ holomorphic in $\D$ with $|\varepsilon_0(z)|\leq 1$, we have:
\begin{align*}
J_\alpha[f](z)-1=\dfrac{1+\varepsilon(z)}{1-\varepsilon(z)}-1 & \Leftrightarrow (1-\alpha)\left(z\dfrac{f'}{f}(z)-1\right)+\alpha z\dfrac{f''}{f'}(z)=\frac{2\varepsilon(z)}{1-\varepsilon(z)}\\
& \Leftrightarrow (1-\alpha)\left(\dfrac{zf'(z)-f(z)}{f(z)}\right)+\alpha z\dfrac{f''}{f'}(z)=\dfrac{2z\varepsilon_0(z)}{1-z\varepsilon_0(z)}.
\end{align*}
Since $f(z)\neq 0$ for $z\neq 0$, which is guarantee by the univalence and normalization, we defined $(zf'(z)-f(z))/f(z)=zh(z)$. Then
    \begin{equation*}
        (1-\alpha)zh(z)+\alpha z\dfrac{f''}{f'}(z)=\dfrac{2z\varepsilon_0(z)}{1-z\varepsilon_0(z)}\Leftrightarrow (1-\alpha) h(z)+\alpha\dfrac{f''}{f'}(z)=\dfrac{2\varepsilon_0(z)}{1-z\varepsilon_0(z)}.
    \end{equation*}
    Let $E=E(\alpha,z)=(1-\alpha)h(z)+\alpha (f''/f')(z)$. Then:
    \begin{equation}\label{eq:phi}
      \varepsilon_0(z)=\dfrac{E}{2+zE}.
    \end{equation}
    Since we know that $\lvert\varepsilon_0(z)\rvert\leq 1$, applying this to \eqref{eq:phi}, we obtain:
    \begin{align*}
        \left\lvert\dfrac{E}{2+zE}\right\rvert\leq 1 & \Leftrightarrow \lvert E\rvert^2\leq\lvert 2+zE\rvert^2\\
        & \Leftrightarrow \lvert E\rvert^2\leq 4+\lvert z\rvert^2\cdot\rvert E\lvert^2+4\text{Re}\{zE\}\\
        & \Leftrightarrow (1-\lvert z\rvert^2)\lvert E\rvert^2\leq 4+4\text{Re}\{zE\}.
    \end{align*}
    Notice that $zE=J_\alpha[f](z)-1$, so we have:
    \begin{equation*}
        (1-\lvert z\rvert^2)\lvert E\rvert^2\leq 4+4\text{Re}\{J_\alpha[f](z)-1\}\Leftrightarrow (1-\lvert z\rvert^2)\lvert E\rvert^2\leq 4\text{Re}\{J_\alpha[f](z)\}.
    \end{equation*}
    Finally, we can conclude:
    \begin{equation*}
        \text{Re}\{J_\alpha[f](z)\}\geq\dfrac{1}{4}(1-\lvert z\rvert^2)\cdot\lvert E\rvert^2.
    \end{equation*}
\end{proof}

\subsection{Fekete--Szeg\H o Inequality} The Fekete--Szeg\H o inequality \cite{FS1933} has been extensively studied in geometric function theory. For univalent mappings $f(z)=z+a_2z^2+a_3z^3+\cdots$ it gives $|a_3-a_2^2|\leq 1$, while for convex mappings ($M_1=C$) the classical sharp bound due by S. Trimble in \cite{Tri1975} shows that $|a_3-a_2^2|\leq (1-|a_2|^2)/3$ holds. The following results generalize this inequality to $\alpha$-convex class.

\begin{theo} Let $\alpha>0$ and $f(z)=z+a_2z^2+a_3z^3+\cdots\in M_\alpha$, then $$|a_3-a_2^2|\leq \dfrac{1}{1+2\alpha}\left[1+\left(\dfrac{|1-\alpha^2|-|1+\alpha|^2}{4}\right)|a_2|^2\right].$$ The bound is sharp and holds when, for $0\leq \lambda\leq 1$, $g(z)=\displaystyle \int_0^z\left(\dfrac{1+t}{1-t}\right)^\lambda \dfrac{1}{1-t^2}dt$ in the integral representation given by \eqref{rep_int}.
\end{theo}

\begin{proof}Since $f$ is starlike, we have that $z\dfrac{f'}{f}(z)=\dfrac{1+\omega}{1-\omega}$ with $\omega:\D\to \D$ holomorphic in $\D$ and $\omega(0)=0$, hence $\omega(z)=z\omega_0(z)$. This implies
\begin{equation}\label{preSch}
\dfrac{f''}{f'}=\dfrac{2\omega'}{1-\omega^2}+\dfrac{2\omega_0}{1-\omega}.\end{equation}
However, $f \in M_\alpha$, therefore there exists $\varepsilon:\D\to \D$ holomorphic in $\D$ with $\varepsilon(0)=0$ such that $J_\alpha[f](z)=(1+\varepsilon(z))/(1-\varepsilon(z))$, from which we have that
\begin{equation}\label{omega_epsilon}
  \dfrac{1+\omega}{1-\omega}+\alpha \dfrac{2z\omega'}{1-\omega^2}=\dfrac{1+\varepsilon}{1-\varepsilon},
\end{equation}
Since $\varepsilon(0)=0$, its follows that $\varepsilon(z)=z\varepsilon_0(z)$. Thus, we have that
\begin{equation*} \frac{\omega_0}{1-\omega}+\frac{\alpha\omega'}{1-\omega^2}=\frac{\varepsilon_0}{1-\varepsilon},\end{equation*} and using equation \eqref{omega_epsilon} we have
\begin{eqnarray}\label{alpha-omega}
\dfrac{\alpha \omega'}{1-\omega^2} & = & \dfrac{\varepsilon_0-\omega_0}{(1-\varepsilon)(1-\omega)}.\label{ec1}
\end{eqnarray}
But $\omega'(z)=\omega_0(z)+z\omega_0'(z)$ and $\omega''(z)=2\omega_0'(z)+z\omega_0''(z)$. Thus, using equation \eqref{preSch}, we have that
$$a_3-a_2^2=\frac{1}{2}\omega''(0)-\omega'(0)^2,$$ and by \eqref{ec1} we have that $\alpha\omega'(0)=\varepsilon'(0)-\omega'(0)$ and from this,
$$\omega'(0)=\dfrac{\varepsilon'(0)}{1+\alpha}.$$
Additionally, differentiating \eqref{alpha-omega}, we obtain
$$\alpha\dfrac{\omega''(1-\omega^2)+2\omega\omega'^2}{(1-\omega^2)^2}=\dfrac{(\varepsilon_0'-\omega_0')(1-\varepsilon)(1-\omega)+(\varepsilon_0-\omega_0)\left(\varepsilon'(1-\omega)+\omega'(1-\varepsilon)\right)}{(1-\varepsilon)^2(1-\omega)^2},$$
evaluating at $z=0$ yields
\begin{eqnarray}
\alpha\omega''(0) & = & \varepsilon'(0)-\omega'(0)+(\varepsilon_0-\omega_0)(\varepsilon'(0)+\omega'(0)).\label{ec2}
\end{eqnarray}

But, $\varepsilon''(0)=2\varepsilon_0'(0)$ and $\omega''(0)=2\omega_0'(0)$. Using these relationships in \eqref{ec2}, we get
$$\omega''(0)=\dfrac{\varepsilon''(0)}{1+2\alpha}+\dfrac{2\alpha(2+\alpha)\varepsilon'(0)^2}{(1+\alpha)^2(1+2\alpha)}.$$

Substituting, we have that:
\begin{eqnarray*}
a_3-a_2^2 & = & \frac{1}{2}\left[\dfrac{\varepsilon''(0)}{1+2\alpha}+\dfrac{2\alpha(2+\alpha)\varepsilon'(0)^2}{(1+\alpha)^2(1+2\alpha)}-\dfrac{2\varepsilon'(0)^2}{(1+\alpha)^2}\right]\\
& = & \dfrac{1}{2(1+2\alpha)}\left[\varepsilon''(0)-\dfrac{2(1-\alpha^2)\varepsilon'(0)^2}{(1+\alpha)^2}\right].
\end{eqnarray*}

If we write $\varepsilon(z)=\varepsilon_1z+\varepsilon_2z^2+\cdots=z(\varepsilon_1+\varepsilon_2z+\cdots)$ and using that $|\varepsilon_2|\leq1-|\varepsilon_1|^2$, we have that $|\varepsilon''(0)|\leq 2-2|\varepsilon'(0)|^2$, and with this
\begin{eqnarray*}
|a_3-a_2^2| & \leq & \dfrac{1}{2(1+2\alpha)}\left[2-2|\varepsilon'(0)|^2+\dfrac{2|1-\alpha^2||\varepsilon'(0)|^2}{(1+\alpha)^2}\right]\\
& = & \dfrac{1}{1+2\alpha}\left[1+\left(\left|\dfrac{1-\alpha}{1+\alpha}\right|-1\right)|\varepsilon'(0)|^2\right].
\end{eqnarray*}
Since $\varepsilon'(0)=(1+\alpha)\omega'(0)$ and $f''(0)=4\omega'(0)$ (see Eq.~\eqref{preSch}), we conclude that
$$|a_3-a_2^2|\leq \dfrac{1}{1+2\alpha}\left[1+\left(\dfrac{|1-\alpha^2|-|1+\alpha|^2}{4}\right)|a_2|^2\right].$$ The equality holds if we consider that $$\varepsilon_0(z)=\frac{z+\lambda}{1+\lambda z},\quad 0\leq \lambda\leq 1,$$ which implies that, in the integral representation \eqref{rep_int} the corresponding function $g$ is given by $$g'(z)=\left(\dfrac{1+z}{1-z}\right)^\lambda \cdot \dfrac{1}{1-z^2}$$
\end{proof}

\subsection{The Order of a function.} On the other hand, the \textit{order} of a mapping \( f \) is defined as the supremum of the modulus of the second Taylor coefficient of the function:
\[
\frac{f(\sigma_a(z)) - f(a)}{f'(a)(1 - |a|^2)},
\]
where \( \sigma_a(z) = (z + a)/(1 + \overline{a}z) \) and \( a \in \D \) . This invariant is denoted by \( A_f \) and is given by
\begin{equation*}\label{Af} A_f=\sup_{z\in\mathbb{D}}\left|\frac{(1-|z|^2)}{2}\frac{f''}{f'}(z)-\overline{z}\right|.\end{equation*}
It is not hard to see that $A_f\geq 1$ for any $f$ and $A_f=1$ if and only if $f$ is a convex mappings. Moreover, if $f$ is univalent, then $A_f\leq 2$. The reader can see Pommerencke's classical paper \cite{P1964}.

\begin{theo} Let $f$ in $M_\alpha$ with $\alpha\in [0,1]$, then its order satisfies that $A_f\leq 2-\alpha.$
\end{theo}

\begin{proof} Subtracting $1$ from both sides of equation \eqref{omega_epsilon} we obtain
\begin{equation*}
    \dfrac{2\omega}{1-\omega}+\dfrac{2\alpha z\omega'}{1-\omega^2}=\dfrac{2\varepsilon}{1-\varepsilon}\,\Leftrightarrow\,\dfrac{2\omega_0}{1-\omega}=\dfrac{2\varepsilon_0}{1-\varepsilon}-\dfrac{2\alpha\omega'}{1-\omega^2}.
\end{equation*}
Replacing the right side of the last equation in equation \eqref{preSch}, we have
\begin{equation*}
    \dfrac{f''}{f'}=\dfrac{2\omega'}{1-\omega^2}+\dfrac{2\varepsilon_0}{1-\varepsilon}-\dfrac{2\alpha\omega'}{1-\omega^2}=\dfrac{2\varepsilon_0}{1-\varepsilon}+\dfrac{2(1-\alpha)\omega'}{1-\omega^2}.
\end{equation*}
Then (here $\varepsilon_0$, $\omega$, and $\omega'$ are evaluating at $z$)
\begin{align*}
\dfrac{1}{2}(1-|z|^2)\dfrac{f''}{f'}(z)-\overline{z} & =  (1-|z|^2)\left(\dfrac{\varepsilon_0}{1-\varepsilon}+(1-\alpha)\dfrac{\omega'}{1-\omega^2}\right)-\overline{z}\\
& =  \dfrac{\varepsilon_0-\overline{z}}{1-z\varepsilon_0}+(1-\alpha)\dfrac{\omega'(1-|z|^2)}{1-\omega^2}.
\end{align*}
Applying Schwarz-Pick lemma and the fact that $|(\varepsilon_0-\overline{z})/(1-z\varepsilon_0)|<1$ for all $|z|<1$ we obtain
\begin{align*}
    \left|\dfrac{1}{2}(1-|z|^2)\dfrac{f''}{f}(z)-\overline{z}\right| & \leq \left|\dfrac{\varepsilon_0-\overline{z}}{1-z\varepsilon_0}\right|+(1-\alpha)\left|\dfrac{\omega'(1-|z|^2)}{1-\omega^2}\right|\\
    & \leq 1+(1-\alpha)(1-|z|^2)\dfrac{|\omega'|}{1-|\omega|^2}\\
    & \leq 1+(1-\alpha)\dfrac{1-|z|^2}{1-|\omega|^2}\cdot\dfrac{1-|\omega|^2}{1-|z|^2}=1+(1-\alpha)=2-\alpha.
\end{align*}
Thus, taking supremum, we have that $A_f(z)\leq 2-\alpha$, which completes the proof.
\end{proof}

\subsection{Schwarzian Derivative} The Schwarzian derivative of a locally univalent mapping \( f \) is
\[S_f(z) = \left( \frac{f''(z)}{f'(z)} \right)' - \frac{1}{2} \left( \frac{f''(z)}{f'(z)} \right)^2.\]
This is one of the most thoroughly studied invariant in complex analysis; see \cite{CDO2011} and the references therein. Since every function in $M_\alpha$ is univalent, its Schwarzian norm satisfies
\[\|Sf\| = \sup_{z \in \mathbb{D}} |Sf(z)|(1-|z|^2)^2 \leq 6.\]
We seek a sharper bound in terms of $\alpha$. As we mentioned in the introduction, if $\alpha \geq 1$, then $M_\alpha \subset S^\ast_\delta$ and $M_\alpha\subset C_\beta$, where $\beta=(1-1/\alpha)\delta(\alpha)$ and $\delta(\alpha)$ is the increasing function given by
\begin{equation}\label{delta}\delta(\alpha):=\frac{\Gamma\left(\frac12+\frac1\alpha\right)}{\sqrt{\pi}\cdot \Gamma(1+\frac1\alpha)}.\end{equation}
It is not hard to see that $0\leq \delta(\alpha)\leq 1$ for all $\alpha\geq1$.

\begin{theo} Let $\alpha\geq 0$ and $f\in M_\alpha$, then $\|Sf\|\leq s(\alpha)$, where
\begin{equation}\label{def-s} s(\alpha)=\left\{\begin{array}{lr}6, & 0\leq \alpha\leq \alpha_0\\[0.3cm]
\dfrac2\alpha\left(2-\alpha^2\right), & \alpha_0<\alpha<1.\\[0.3cm]
2, & 1\leq \alpha\leq \alpha_1.\\[0.3cm]
8\left(1-\frac1\alpha\right)\delta(\alpha)\left(1-\left(1-\frac1\alpha\right)\delta(\alpha)\right), & \alpha>\alpha_1.\\[0.3cm]
\end{array}\right.\end{equation}
Here $\alpha_0=(\sqrt{17}-3)/2\sim 0.56$ and $\alpha_1$ is the number that satisfies $\left(1-\frac{1}{\alpha_1}\right)\delta(\alpha_1)=1/2$ ($\alpha_1\sim 3.27$).
\end{theo}
\begin{proof} For $\alpha\in[0,1]$ we have that $f\in S^\ast$ then $zf'(z)/f(z)=(1+\omega(z))/(1-\omega(z))$ and
\begin{equation}\label{conv-omega}1+z\frac{f''}{f'}(z)=\frac{z\omega'(z)}{1+\omega(z)}+\frac{z\omega'(z)}{1-\omega(z)}+\frac{1+\omega(z)}{1-\omega(z)}.\end{equation}
Also $J_\alpha[f](z)=(1+\varepsilon(z))/(1-\varepsilon(z))$, therefore
\begin{equation}\label{conv-eps}1+z\frac{f''}{f'}(z)=\frac1\alpha\cdot \frac{1+\varepsilon(z)}{1-\varepsilon(z)}-\frac{1-\alpha}{\alpha}\cdot\frac{1+\omega(z)}{1-\omega(z)}.\end{equation}
Recall that $\varepsilon(z)=z\varepsilon_0(z)$ and $\omega(z)=z\omega_0(z)$ as before. Thus, (we omitted the argument $(z)$), $$\frac{f''}{f'}=\frac2\alpha\left(\frac{\varepsilon_0}{1-\varepsilon}-\frac{(1-\alpha)\omega_0}{1-\omega}\right).$$ Then,
$$ Sf=\frac{2}{\alpha}\left(\frac{\varepsilon_0'+\varepsilon_0^2}{(1-\varepsilon)^2}-(1-\alpha)\frac{\omega_0'+\omega_0^2}{(1-\omega)^2}\right)-\dfrac{2}{\alpha^2}\left(\dfrac{\varepsilon_0}{1-\varepsilon}-\dfrac{(1-\alpha)\omega_0}{1-\omega}\right)^2,$$
which is equivalent to $$Sf=\frac{2}{\alpha}\left(\frac{\varepsilon_0'}{(1-\varepsilon)^2}-(1-\alpha)\frac{\omega_0'}{(1-\omega)^2}\right)-\dfrac{2(1-\alpha)}{\alpha^2}\left(\dfrac{\varepsilon_0}{1-\varepsilon}-\dfrac{\omega_0}{1-\omega}\right)^2.$$
Using equation \eqref{omega_epsilon}, we have that $$Sf=\frac{2}{\alpha}\left(\frac{\varepsilon_0'}{(1-\varepsilon)^2}-(1-\alpha)\frac{\omega_0'}{(1-\omega)^2}\right)-2(1-\alpha)\left(\dfrac{\omega'}{1-\omega^2}\right)^2.$$
Using Schwarz-Pick Lemma and $1-|\varphi_g(z)|^2=(1-|g(z)|^2)(1-|z|^2)/|1-zg(z)|^2$, where $\varphi_g(z)=(\overline z-g(z))/(1-zg(z))$, we have that
\begin{equation}\label{bound-sf}\begin{array}{lll}
 |Sf(z)|(1-|z|^2)^2 &\leq & \dfrac{2}{\alpha}\left(1-|\varphi_{\varepsilon_0}(z)|^2+(1-\alpha)(1-|\varphi_{\omega_0}(z)|^2)\right)+2(1-\alpha)\left(\dfrac{|\omega'|(1-|z|^2)}{1-|\omega|^2}\right)^2\\[0.3cm]
  &\leq  & \dfrac{2}{\alpha}\left(1+1-\alpha\right)+2(1-\alpha)= \dfrac{2}{\alpha}\left(2-\alpha^2\right).
\end{array} \end{equation}
Note that $|Sf(z)|(1-|z|^2)^2 \leq 6$ whenever $\alpha \geq (\sqrt{17}-3)/2 \approx 0.56$, then the first two part of the $s(\alpha)$ has been proved. On the other hand, when $\alpha\geq 1$, we have that $f\in C_\beta$, given by $\beta=(1-1/\alpha)\delta(\alpha)$. In fact, Suita in \cite{Sui1996} proved that, for any $f\in C_\beta$, then the norm of the Schwarzian derivative is bounded by 2 when $0\leq \beta\leq 1/2$, and $8\beta(1-\beta)=8\left(1-\frac1\alpha\right)\delta(\alpha)\left(1-\left(1-\frac1\alpha\right)\delta(\alpha)\right)$ when $1/2<\beta<1$. As $\beta=1/2$ occur when $\left(1-\frac{1}{\alpha_0}\right)\delta(\alpha_0)=1/2$, the proof is completed.
\end{proof}

\begin{coro} For $0<\alpha<1$, we have that
\begin{equation*}|Sf(z)|(1-|z|^2)^2+ \frac{2}{\alpha}\left|\frac{\overline z-\varepsilon_0(z)}{1-\varepsilon(z)}\right|^2+(1-\alpha)\left|\frac{\overline z-\omega_0(z)}{1-\omega(z)}\right|^2\leq \dfrac{2}{\alpha}(2-\alpha^2).\end{equation*}
\end{coro}

\begin{proof} From the proof of the previous Theorem, using inequality \eqref{bound-sf} we can rearrange to obtain this.
\end{proof}

\subsubsection{The Schwarzian derivative of $\alpha$-Koebe function}

We end this work by showing an example associated with the integral representation of any given function in $M_\alpha$ for $\alpha>0$. In \cite{MMR1974} the authors shows that $f\in M_\alpha$ ($\alpha>0$) if and only if there exists $g\in S^*$ such that
\begin{equation}
    f(z)=\left[\dfrac{1}{\alpha}\int_0^z\dfrac{(g(\xi))^{1/\alpha}}{\xi}\,d\xi\right]^\alpha.
\end{equation}
This integral representation allows us to construct $\alpha$-convex functions given any starlike function. In particular, if $k(z)=z/(1-z)^2$ ($z\in \D$) is the Koebe function we obtain the $\alpha$-convex Koebe function
\begin{equation}
    k_\alpha(z)=\left[\dfrac{1}{\alpha}\int_0^z \xi^{1/\alpha-1}(1-\xi)^{-2/\alpha}\,d\xi\right]^{\alpha}.
\end{equation}
It is well known that the Koebe function maximizes many important geometric properties in certain classes of univalent functions; in particular, $k_\alpha$ yields sharp growth estimates in $M_\alpha$ \cite{M1973}.

From an algebraic point of view in equation \eqref{rep_int}, let $G(z)=\int_0^z(g(\xi))^{1/\alpha}\cdot\xi^{-1}\,d\xi$, then $f(z)=[(1/\alpha)\cdot G(z)]^\alpha$ then
\begin{equation*}
    \dfrac{f''}{f'}(z)=(\alpha-1)\dfrac{G'}{G}(z)+\dfrac{G''}{G'}(z).
\end{equation*}
Thus, we can obtain an expression for the Schwarzian derivative in terms of $G$
\begin{equation}\label{schwarz-G}
    Sf(z)=\left(\dfrac{f''}{f'}(z)\right)'-\dfrac{1}{2}\left(\dfrac{f''}{f'}(z)\right)^2=SG(z)-\dfrac{(\alpha-1)^2}{2}\left(\dfrac{G'}{G}(z)\right)^2.
\end{equation}
Notice that $G$ is not necessarily locally univalent in $\D$. Now if we let $n=1/\alpha$ for $n\in\mathbb{N}$ with $n\geq 2$, then \eqref{koebeconvex} can be expressed as
\begin{equation}\label{koebeconvex2}
    k_n(z)=\left[n\int_0^z \xi^{n-1}(1-\xi)^{-2n}\,d\xi\right]^{1/n}=[n\text{B}(z;n,1-2n)]^{1/n},
\end{equation}
where $B(z;a,b)$ is the incomplete beta function. Notice that the function in the integrand is rational so if we integrate, the expression \eqref{koebeconvex2} can also be represented as
\begin{equation*}
    k_n(z)=\left[n\sum_{k=0}^{n-1}(-1)^k\binom{n-1}{k}\left(\dfrac{1-(1-z)^{k-2n+1}}{k-2n+1}\right)\right]^{1/n}.
\end{equation*}
In this case $G(z)=B[z;n,1-2n]$ and apply \eqref{schwarz-G} with $\alpha=1/n$, we obtain
\begin{equation}\label{Schwarz-k}
    Sk_n(z)=\dfrac{1-n}{2z^2}\left[1+n\left(\dfrac{1+z}{1-z}\right)^2\right]+\left(\dfrac{n^2-1}{2n^2}\right)\left(\dfrac{G'}{G}(z)\right)^2.
\end{equation}
It is not hard to see that $Sk_n(0)=6n(n-1)/((n+2)(n+1))$, then $$\sup_{f\in M_\alpha}\|Sf\|\geq \frac{6n(n-1)}{(n+2)(n+1)}.$$

For $\alpha = 1/n$, the expression $Sk_n$ is a rational function of $z$ and $G(z)$. This allows us to compute the Schwarzian norm numerically for each $n$. For $n = 1, \ldots, 9$, a direct computation (using Mathematica) yields
\begin{equation}\label{schwarz-k-n}\|Sk_n\|=\sup_{z\in\D}|Sk_n(z)|(1-|z|^2)^2=2\left(3-\frac1n\right)\left(1-\frac1n\right), \quad \quad n=1,\ldots, 9.\end{equation}
Numerical evidence for $10 \leq n < 50$ is consistent with this formula, as Figure~\ref{fig1} illustrates. Thus, we conjectured that $$\sup_{f\in M_\alpha}\|Sf\|=2\left(3-\frac1n\right)\left(1-\frac1n\right).$$

\begin{figure}[h]
\centering
\includegraphics[width=\textwidth]{Fig1.png}
\caption{(Left) Points in $\D$. (Right) $|Sk_n(z)|(1-|z|^2)^2$.}
\label{fig1}
\end{figure}

\begin{thebibliography}{99}

\bibitem{CDO2011} M. Chuaqui, P. Duren, and B. Osgood, Schwarzian derivatives of convex mappings, Ann. Acad. Sci. Fenn. Math. \textbf{36} (2011), no.~2, 449--460.

\bibitem{FS1933} M. Fekete and G. Szeg\H{o}, Eine Bemerkung \"uber ungerade schlichte Funktionen, J. Lond. Math. Soc. \textbf{8} (1933), no.~4, 85--89.

\bibitem{GK2003} I. Graham and G. Kohr, \emph{Geometric Function Theory in One and Higher Dimensions}, Marcel Dekker, New York, 2003.

\bibitem{M1973} S. Miller, Distortion properties of alpha-starlike functions, Proc. Amer. Math. Soc. \textbf{38} (1973), 311--318.

\bibitem{MMR1973} S. Miller, P. Mocanu, and M. Reade, All $\alpha$-convex functions are univalent and starlike, Proc. Amer. Math. Soc. \textbf{37} (1973), 553--554.

\bibitem{MMR1974} S. Miller, P. Mocanu, and M. Reade, Bazilevi\v{c} functions and generalized convexity, Rev. Roumaine Math. Pures Appl. \textbf{19} (1974), 213--224.

\bibitem{MMR1978} S. Miller, P. Mocanu, and M. Reade, The order of starlikeness of alpha-convex functions, Mathematica (Cluj) \textbf{20} (1978), no.~43, 25--30.

\bibitem{MM2000} S. S. Miller and P. T. Mocanu, \emph{Differential Subordinations: Theory and Applications}, Marcel Dekker, New York, 2000.

\bibitem{Moc1969} P. Mocanu, Une propri\'et\'e de convexit\'e g\'en\'eralis\'ee dans la th\'eorie de la repr\'esentation conforme, Mathematica (Cluj) \textbf{11} (1969), no.~34, 127--133.

\bibitem{P1964} Ch. Pommerenke, Linear-invariante Familien analytischer Funktionen I, Math. Ann. \textbf{155} (1964), 108--154.

\bibitem{Sui1996} N. Suita, Schwarzian derivatives of convex functions, J. Hokkaido Univ. Ed. Sect. II A \textbf{46} (1996), no.~2, 113--117.

\bibitem{Tri1975} S. Y. Trimble, A coefficient inequality for convex univalent functions, Proc. Amer. Math. Soc. \textbf{48} (1975), no.~1, 266--267.

\end{thebibliography}
\end{document}
